% !PRINTED 535
\weberpage{535}

Wie man aus der einen Basis von $\mu$ alle anderen ableiten kann, haben
wir schon oben gesehen.

Wir verstehen jetzt unter $\alpha_1, \alpha_2, \ldots, \alpha_n$ eine
beliebige Basis von $\mu$ und stellen das Functional $\mu$ nach
\S.\,142, 4. als Quotienten zweier ganzer Functionen dar, dessen Nenner
eine rationale Einheit ist. Die Co\"efficienten des Zählers sind dann
ganze durch $\mu$ theilbare Zahlen (\S.\,143, 2.) und können daher in
der Form (2) dargestellt werden.

Fassen wir diese Darstellung gehörig zusammen, so ergiebt sich also für
$\mu$ ein Ausdruck
\weq{8}{
\mu = \alpha_1 u_1 + \alpha_2 u_2 + \cdots + \alpha_n u_n,
}
worin die $u_1, u_2, \ldots, u_n$ \gesperrt{ganze rationale Functionale}
sind.

Hiermit wollen wir die Linearform
\weq{9}{
\lambda = \alpha_1 t_1 + \alpha_2 t_2 + \cdots + \alpha_n t_n
}
vergleichen, in der $t_1, t_2, \ldots, t_n$ Variable sind. Diese
Linearform ist (nach \S.\,139) der grösste gemeinschaftliche Theiler der
Zahlen $\alpha_1, \alpha_2, \ldots, \alpha_n$ und ist durch $\mu$
theilbar, weil die Zahlen $\alpha_1, \alpha_2, \ldots, \alpha_n$ durch
$\mu$ theilbar sind. Andererseits ist aber auch, wie die Darstellung (8)
zeigt, $\mu$ durch $\lambda$ theilbar, und folglich sind die beiden
Functionale $\mu$ und $\lambda$ mit einander associirt.

\gesperrt{Das Functional $\lambda$ wollen wir eine Basisform des
Functionals $\mu$ nennen; es ist dann $\lambda$ zugleich Basisform von
allen mit $\mu$ associirten Functionalen.}

Eine Basisform $\lambda$ von $\mu$ hat die Eigenschaft, dass
man aus ihr alle durch $\mu$ (oder durch $\lambda$) theilbaren ganzen
Zahlen in $\Omega$ erhält, wenn man für die Variablen ganze rationale
Zahlen setzt.

Die Basis $\alpha_1, \alpha_2, \ldots, \alpha_n$ steht zu dem
Functionale $\mu$ in einer ähnlichen Beziehung, wie die Basis $\omega_1,
\omega_2, \ldots, \omega_n$ des Systemes $\mathfrak{o}$ aller ganzen Zahlen in
$\Omega$ zu den Einheiten. In der That ist die Linearform
\weq{10}{
\tau = \omega_1 t_1 + \omega_2 t_2 + \cdots + \omega_n t_n
}
eine Einheit; denn sie ist der grösste gemeinschaftliche Theiler von
$\omega_1, \omega_2, \ldots, \omega_n$ und muss also, wie die Formel
\S.\,145, (8) zeigt, ein Theiler von 1, also eine Einheit sein.

\gesperrt{Demnach wollen wir das Functional $\tau$ eine Basisform von $\mathfrak{o}$
nennen.}
